\documentclass[conference]{IEEEtran}

% Packages
\usepackage[utf8]{inputenc}
\usepackage[T1]{fontenc}
\usepackage{graphicx}
\usepackage{booktabs}
\usepackage{hyperref}
\usepackage{amsmath}
\usepackage{algorithm}
\usepackage{algpseudocode}
\usepackage{caption}
\usepackage{subcaption}
\usepackage{float}
\usepackage{enumitem}

% hyperref setup
\hypersetup{
    colorlinks=true,
    linkcolor=blue,
    citecolor=blue,
    urlcolor=blue
}

\title{Bangladesh Suicide Risk Detection: A Cross-Cultural Analysis Using News Reports}

\author{
\IEEEauthorblockN{Author One}
\IEEEauthorblockA{Department of Computer Science\\University Name\\Bangladesh\\author1@email.com}
\and
\IEEEauthorblockN{Author Two}
\IEEEauthorblockA{Department of Computer Science\\University Name\\Bangladesh\\author2@email.com}
\and
\IEEEauthorblockN{Author Three}
\IEEEauthorblockA{Department of Computer Science\\University Name\\Bangladesh\\author3@email.com}
}

\begin{document}

\maketitle

\begin{abstract}
This paper presents a comprehensive approach to suicide risk detection in Bangladesh using newspaper-reported incidents. We propose a multi-stage pipeline that combines statistical imputation, cross-cultural feature engineering, and clustering algorithms to identify regional patterns and risk factors. Our method leverages multilingual LLM inference and cross-cultural feature engineering to extract meaningful insights from Bangla-language news reports. The results demonstrate the effectiveness of our approach in identifying key suicide risk factors across different regions of Bangladesh, providing valuable insights for public health policy and NGO interventions.
\end{abstract}

\begin{IEEEkeywords}
Suicide risk detection, Bangladesh, Natural Language Processing, Clustering, Public Health
\end{IEEEkeywords}

%----------------------------------------------------------
\section{Introduction}
%----------------------------------------------------------

Suicide is a major public health concern globally, with significant variations across different regions and cultures. Bangladesh, being a densely populated South Asian country, faces unique challenges related to mental health and suicide prevention.

The lack of comprehensive datasets and culturally-aware analytical methods makes it challenging to understand suicide risk factors in the Bangladeshi context. Existing studies often fail to account for regional, cultural, and socioeconomic variations that influence suicidal behavior.

\textbf{Objectives:}
\begin{itemize}
    \item Develop a systematic approach to extract suicide-related information from Bangla news reports
    \item Identify regional patterns and risk factors through cross-cultural analysis
    \item Propose a scalable framework for suicide risk detection in similar contexts
\end{itemize}

\textbf{Contributions:}
\begin{enumerate}
    \item A comprehensive dataset of suicide incidents extracted from Bangla newspapers
    \item A multi-stage pipeline combining statistical imputation, LLM inference, and clustering
    \item Analysis of regional triggers and demographic patterns
    \item Policy recommendations for targeted interventions
\end{enumerate}

%----------------------------------------------------------
\section{Dataset}
%----------------------------------------------------------

\subsection{Data Collection}
The dataset was constructed by scraping articles from major Bangla newspapers including Prothom Alo. The collection process involved keyword-based search using Bangla suicide-related terms, covering all divisions of Bangladesh.

\subsection{Preprocessing}
The preprocessing pipeline includes:
\begin{enumerate}
    \item Language detection and translation (Bangla to English)
    \item Text normalization and deduplication
    \item Named entity anonymization (case\_id hashing)
    \item Geographic coordinate extraction (district centroids)
\end{enumerate}

%----------------------------------------------------------
\section{Method}
%----------------------------------------------------------

Figure~\ref{fig:pipeline} illustrates the complete methodology pipeline.

\begin{figure}[H]
\centering
\fbox{\parbox{0.7\columnwidth}{\centering\vspace{2cm}Pipeline Figure Placeholder\vspace{2cm}}}
\caption{Methodology Pipeline}
\label{fig:pipeline}
\end{figure}

\subsection{Imputation \& Feature Extraction}
We employ statistical imputation methods (MissForest, MICE, KNN) and multilingual LLM inference (Llama-3, IndicBERT, Gemini) for feature extraction. Cross-cultural feature engineering captures regional triggers and demographic patterns.

\subsection{Cross-Regional Clustering}
Text data is converted to dense vector embeddings using transformer-based models. We apply Agglomerative Clustering, K-Means, and DBSCAN algorithms, supplemented by BERTopic for thematic analysis.

\subsection{Risk Scoring \& Validation}
A fine-tuned LLM risk classifier categorizes regional triggers (Farm Distress, Dowry, Academic, Economic). Custom cluster validation metrics ensure quality clustering.

%----------------------------------------------------------
\section{Results}
%----------------------------------------------------------

\subsection{Imputation Performance}

\begin{table}[H]
\centering
\caption{Imputation Method Performance}
\label{tab:imputation}
\begin{tabular}{@{}lccc@{}}
\toprule
\textbf{Method} & \textbf{RMSE} & \textbf{MAE} & \textbf{R$^2$} \\
\midrule
MissForest & -- & -- & -- \\
MICE & -- & -- & -- \\
KNN & -- & -- & -- \\
\bottomrule
\end{tabular}
\end{table}

\subsection{Clustering Results}
The clustering analysis revealed distinct regional patterns across Bangladesh.

\subsection{Risk Factor Analysis}
Key risk factors identified include economic distress, family conflict, academic pressure, dowry-related issues, and farm distress.

%----------------------------------------------------------
\section{Discussion}
%----------------------------------------------------------

The analysis reveals several important patterns in Bangladesh suicide incidents. Our approach aligns with and extends previous studies in suicide risk detection.

\textbf{Limitations:} Data availability constraints, language processing challenges, and generalizability concerns.

\textbf{Future Work:} Expand to additional South Asian countries, incorporate real-time monitoring, and develop intervention frameworks.

%----------------------------------------------------------
\section{Conclusions}
%----------------------------------------------------------

This paper presented a comprehensive approach to suicide risk detection in Bangladesh. Our multi-stage pipeline successfully identifies regional patterns and risk factors from newspaper reports. The findings provide valuable insights for public health policy and targeted interventions.

%----------------------------------------------------------
% References
%----------------------------------------------------------
\begin{thebibliography}{9}

\bibitem{example}
A. Author and B. Author, ``Title of the article,'' \textit{Journal Name}, vol.~1, no.~1, pp.~1--10, 2021.

\bibitem{llm}
J. Smith et al., ``Large language models for text analysis,'' in \textit{Proc. of NLP Conf.}, 2023, pp.~123--145.

\bibitem{clustering}
R. Kumar et al., ``Clustering approaches for document analysis,'' in \textit{Proc. of ACM Computing Surveys}, 2022, pp.~1--30.

\end{thebibliography}

\end{document}
